%% ========================================================================
%% Chapter 1: Pengenalan Sistem Operasi & Instalasi
%% ========================================================================

\chapter{Pengenalan Sistem Operasi \& Instalasi}
\label{ch:introduction}

Bab ini memperkenalkan konsep fundamental sistem operasi, arsitektur kernel, perbandingan sistem operasi modern, serta proses instalasi dan konfigurasi sistem operasi dengan fokus praktis menggunakan Linux.

%% ========================================================================
%% SESSION 1-2: KONSEP DASAR & ARSITEKTUR OS
%% ========================================================================

%% ------------------------------------------------------------------------
\section{Definisi dan Fungsi Sistem Operasi}
\label{sec:os-basics}

Sistem operasi adalah perangkat lunak sistem yang mengelola sumber daya perangkat keras komputer dan menyediakan layanan umum untuk program aplikasi. Sistem operasi bertindak sebagai perantara antara pengguna dengan perangkat keras komputer.

\subsection{Peran Ganda Sistem Operasi}

Sistem operasi memiliki dua peran utama:

\begin{enumerate}
    \item \textbf{Resource Manager}: Mengelola dan mengalokasikan sumber daya sistem (CPU, memory, disk, I/O devices) secara efisien dan adil antar berbagai program yang berjalan.

    \item \textbf{Extended Machine}: Menyediakan abstraksi dari hardware yang kompleks, memberikan interface yang lebih mudah digunakan untuk programmer dan pengguna.
\end{enumerate}

\begin{notebox}
Tanpa sistem operasi, setiap program harus memiliki driver sendiri untuk setiap perangkat keras, yang sangat tidak efisien dan sulit dipelihara.
\end{notebox}

\subsection{Fungsi Utama Sistem Operasi}

Sistem operasi modern memiliki lima fungsi utama:

\subsubsection{Process Management}

Sistem operasi bertanggung jawab untuk:
\begin{itemize}
    \item \textbf{Process scheduling}: Menentukan proses mana yang mendapat CPU time
    \item \textbf{Process creation dan termination}: Membuat dan mengakhiri proses
    \item \textbf{Process synchronization}: Mengkoordinasikan multiple processes
    \item \textbf{Inter-process communication (IPC)}: Memfasilitasi komunikasi antar proses
\end{itemize}

Contoh implementasi:
\begin{itemize}
    \item \textbf{Windows}: Task Manager menampilkan proses yang berjalan dan penggunaan sumber daya
    \item \textbf{Linux}: Perintah \cmd{ps}, \cmd{top}, dan \cmd{htop} untuk memantau proses
    \item \textbf{macOS}: Activity Monitor untuk mengelola proses
\end{itemize}

\subsubsection{Memory Management}

Fungsi memory management meliputi:
\begin{itemize}
    \item \textbf{Memory allocation}: Memberikan memori ke proses yang membutuhkan
    \item \textbf{Virtual memory}: Menggunakan disk sebagai extension dari RAM
    \item \textbf{Memory protection}: Mencegah satu proses mengakses memori proses lain
    \item \textbf{Paging dan swapping}: Teknik untuk optimasi penggunaan memori
\end{itemize}

\begin{examplebox}[Virtual Memory]
Ketika RAM penuh, sistem operasi memindahkan sebagian data ke disk (swap space di Linux, pagefile di Windows) sehingga aplikasi dapat terus berjalan, meskipun dengan kinerja yang lebih lambat.
\end{examplebox}

\subsubsection{File Management}

Sistem operasi menyediakan:
\begin{itemize}
    \item \textbf{Organisasi file system}: Struktur hirarkis (direktori, subdirektori)
    \item \textbf{Operasi file}: Membuat, membaca, menulis, menghapus file
    \item \textbf{Kontrol akses}: Izin (permissions) untuk keamanan
    \item \textbf{Mounting file system}: Mengintegrasikan perangkat penyimpanan ke sistem
\end{itemize}

Contoh file systems:
\begin{itemize}
    \item Windows: NTFS (New Technology File System)
    \item macOS: APFS (Apple File System)
    \item Linux: ext4, XFS, Btrfs
\end{itemize}

\subsubsection{I/O Management}

Manajemen Input/Output mencakup:
\begin{itemize}
    \item \textbf{Device drivers}: Perangkat lunak untuk berkomunikasi dengan perangkat keras
    \item \textbf{Buffering}: Penyimpanan sementara untuk kelancaran operasi I/O
    \item \textbf{Interrupt handling}: Merespon sinyal dari perangkat keras
    \item \textbf{Spooling}: Antrean untuk perangkat seperti printer
\end{itemize}

\subsubsection{Security dan Protection}

Aspek keamanan meliputi:
\begin{itemize}
    \item \textbf{Authentication}: Verifikasi identitas pengguna (password, biometrik)
    \item \textbf{Authorization}: Kontrol akses ke sumber daya berdasarkan izin (permissions)
    \item \textbf{Encryption}: Proteksi data (BitLocker di Windows, FileVault di macOS)
    \item \textbf{Auditing}: Pencatatan aktivitas sistem untuk pemantauan keamanan
\end{itemize}

%% ------------------------------------------------------------------------
\section{Jenis dan Kategori Sistem Operasi}
\label{sec:os-categories}

Sistem operasi dapat dikategorikan berdasarkan berbagai kriteria:

\subsection{Berdasarkan Platform}

\begin{description}
    \item[Desktop OS] Dirancang untuk komputer pribadi dengan antarmuka pengguna yang ramah. Contoh: Windows 10/11, macOS, Ubuntu Desktop, Linux Mint.

    \item[Server OS] Dioptimasi untuk menjalankan layanan dan melayani banyak klien. Biasanya tanpa GUI atau dengan GUI minimal. Contoh: Windows Server, Ubuntu Server, Red Hat Enterprise Linux, FreeBSD.

    \item[Mobile OS] Dirancang untuk smartphone dan tablet dengan antarmuka sentuh dan efisiensi daya. Contoh: Android, iOS, HarmonyOS.

    \item[Real-Time OS (RTOS)] Untuk sistem tertanam (embedded systems) dengan batasan waktu yang ketat. Digunakan di otomotif, kontrol industri, perangkat medis. Contoh: FreeRTOS, VxWorks, QNX.

    \item[Embedded OS] Untuk perangkat dengan sumber daya terbatas (IoT, router, perangkat pintar). Contoh: Embedded Linux, Windows IoT.
\end{description}

\subsection{Perbandingan Sistem Operasi Populer}

Tabel~\ref{tab:os-comparison} menunjukkan perbandingan sistem operasi modern yang paling banyak digunakan.

\begin{table}[htbp]
    \centering
    \caption{Perbandingan Sistem Operasi Modern}
    \label{tab:os-comparison}
    \begin{tblr}{
        colspec = {lX[l]X[l]X[l]X[l]},
        width = \linewidth,
        hlines,
        vlines,
        row{1} = {font=\bfseries},
        cells = {font=\small}
    }
        \textbf{OS} & \textbf{Kernel Type} & \textbf{Primary Use} & \textbf{License} & \textbf{Market Share} \\
        Windows & Hybrid & Desktop/Business & Proprietary & Desktop: 75\% \\
        macOS & Hybrid (XNU) & Creative/Dev & Proprietary & Desktop: 15\% \\
        Linux & Monolithic & Server/Dev & Open Source & Server: 60\% \\
        Android & Linux-based & Mobile & Open Source & Mobile: 70\% \\
        iOS & Hybrid (XNU) & Mobile & Proprietary & Mobile: 28\% \\
    \end{tblr}
\end{table}

\begin{tipbox}
Pemilihan sistem operasi bergantung pada kebutuhan:
\begin{itemize}
    \item \textbf{Gaming \& Aplikasi Bisnis}: Windows (kompatibilitas terbaik)
    \item \textbf{Pekerjaan Kreatif}: macOS (standar industri untuk video/audio editing)
    \item \textbf{Server \& Cloud}: Linux (stabilitas, keamanan, tanpa biaya lisensi)
    \item \textbf{Development}: macOS atau Linux (tools berbasis Unix)
\end{itemize}
\end{tipbox}

%% ------------------------------------------------------------------------
\section{Arsitektur Kernel Sistem Operasi}
\label{sec:kernel-architecture}

Kernel adalah inti dari sistem operasi yang memiliki kontrol penuh atas seluruh sistem. Kernel menyediakan abstraksi antara hardware dan software.

\subsection{Apa itu Kernel?}

\textbf{Kernel} adalah program yang selalu berjalan (resident) di memori dan memiliki akses eksklusif ke perangkat keras. Kernel bertanggung jawab untuk:

\begin{itemize}
    \item Manajemen proses tingkat rendah (low-level)
    \item Manajemen memori
    \item Device drivers
    \item System calls
    \item Penanganan interrupt
\end{itemize}

\subsection{Kernel Space vs User Space}

Sistem operasi modern memisahkan memori menjadi dua area:

\begin{description}
    \item[Kernel Space (Ring 0)] Area memori dengan privilege tertinggi. Hanya kode kernel yang berjalan di sini. Akses penuh ke semua perangkat keras dan memori.

    \item[User Space (Ring 3)] Area memori untuk aplikasi pengguna. Akses terbatas ke perangkat keras. Harus melalui system calls untuk mengakses layanan kernel.
\end{description}

\begin{figure}[htbp]
    \centering
    \begin{tikzpicture}[node distance=1.8cm, scale=0.95]
        % User Space
        \node[draw, rectangle, minimum width=11cm, minimum height=1.5cm, fill=blue!10, align=center] (user) {
            \textbf{User Space (Ring 3)}\\[0.2em]
            \small Applications: Browser, Text Editor, Media Player
        };

        % System Call Interface
        \node[draw, rectangle, minimum width=11cm, minimum height=0.8cm, below=1cm of user, fill=yellow!20] (syscall) {\textbf{System Call Interface}};

        % Kernel Space
        \node[draw, rectangle, minimum width=11cm, minimum height=2cm, below=1cm of syscall, fill=red!10, align=center] (kernel) {
            \textbf{Kernel Space (Ring 0)}\\[0.2em]
            \small Process Scheduler, Memory Manager, File Systems, Device Drivers
        };

        % Hardware
        \node[draw, rectangle, minimum width=11cm, minimum height=1.2cm, below=1cm of kernel, fill=gray!20, align=center] (hw) {
            \textbf{Hardware}\\[0.2em]
            \small CPU, RAM, Disk, Network, GPU
        };

        % Arrows
        \draw[->, thick] (user.south) -- (syscall.north) node[midway, right, font=\tiny] {System Calls};
        \draw[->, thick] (syscall.south) -- (kernel.north);
        \draw[->, thick] (kernel.south) -- (hw.north) node[midway, right, font=\tiny] {Direct Access};
    \end{tikzpicture}
    \caption{Arsitektur Layer Sistem Operasi}
    \label{fig:os-layers}
\end{figure}

\begin{notebox}
\textbf{Protection Rings}: x86 architecture memiliki 4 privilege levels (Ring 0-3), tapi modern OS hanya menggunakan Ring 0 (kernel) dan Ring 3 (user). Ring 1 dan 2 jarang digunakan.
\end{notebox}

\subsection{Jenis Arsitektur Kernel}

\subsubsection{Monolithic Kernel}

\textbf{Karakteristik:}
\begin{itemize}
    \item Semua layanan (file system, device drivers, manajemen memori) berjalan di kernel space
    \item Binary tunggal yang besar
    \item Pemanggilan fungsi langsung antar komponen (cepat)
\end{itemize}

\textbf{Keuntungan:}
\begin{itemize}
    \item \textbf{Kinerja tinggi}: Tidak ada overhead untuk context switching
    \item \textbf{Komunikasi efisien}: Pemanggilan fungsi lebih cepat dari IPC
    \item \textbf{Desain lebih sederhana}: Semua dalam satu address space
\end{itemize}

\textbf{Kerugian:}
\begin{itemize}
    \item \textbf{Keandalan}: Bug di satu driver dapat menyebabkan crash seluruh sistem
    \item \textbf{Keamanan}: Permukaan serangan lebih luas
    \item \textbf{Pemeliharaan}: Sulit untuk mengisolasi dan debug
\end{itemize}

\textbf{Contoh:} Linux, traditional Unix, BSD

\subsubsection{Microkernel}

\textbf{Karakteristik:}
\begin{itemize}
    \item Kernel minimal dengan hanya layanan esensial (IPC, penjadwalan dasar, memori tingkat rendah)
    \item Sebagian besar layanan (file systems, drivers, network stack) di user space
    \item Message passing untuk komunikasi antar komponen
\end{itemize}

\textbf{Keuntungan:}
\begin{itemize}
    \item \textbf{Keandalan}: Crash layanan tidak mempengaruhi kernel
    \item \textbf{Keamanan}: Permukaan serangan lebih kecil, isolasi lebih baik
    \item \textbf{Modular}: Mudah menambah/menghapus layanan
    \item \textbf{Cocok untuk real-time}: Perilaku yang dapat diprediksi
\end{itemize}

\textbf{Kerugian:}
\begin{itemize}
    \item \textbf{Overhead kinerja}: Context switches dan IPC lebih lambat
    \item \textbf{Kompleksitas}: Message passing menambah kompleksitas
\end{itemize}

\textbf{Contoh:} MINIX 3, QNX, seL4, L4 family

\subsubsection{Hybrid Kernel}

\textbf{Karakteristik:}
\begin{itemize}
    \item Mengkombinasikan ide dari monolithic dan microkernel
    \item Beberapa layanan di kernel space untuk kinerja
    \item Beberapa layanan di user space untuk modularitas
    \item Keseimbangan antara kecepatan dan keandalan
\end{itemize}

\textbf{Contoh:}
\begin{itemize}
    \item \textbf{Windows NT kernel}: Digunakan di Windows, Xbox
    \item \textbf{macOS XNU (X is Not Unix)}: Mach microkernel + BSD components
\end{itemize}

\begin{figure}[htbp]
    \centering
    \resizebox{0.85\textwidth}{!}{%
    \begin{tikzpicture}
        % Monolithic
        \node[draw, rectangle, minimum width=3.5cm, minimum height=5cm] (mono) at (0,0) {};
        \node[above] at (mono.north) {\textbf{Monolithic}};
        \node[draw, fill=red!20, minimum width=3cm, minimum height=4.5cm] at (mono.center) {};
        \node[text width=2.5cm, align=center] at (mono.center) {Kernel\\[0.5em]All Services};
        \node[below=0.1cm of mono.south, text width=3.5cm, align=center, font=\small] {Linux, BSD};

        % Microkernel
        \node[draw, rectangle, minimum width=3.5cm, minimum height=5cm] (micro) at (5,0) {};
        \node[above] at (micro.north) {\textbf{Microkernel}};
        \node[draw, fill=red!20, minimum width=3cm, minimum height=1cm] at (micro.center) {\small Minimal Kernel};
        \node[draw, fill=blue!20, minimum width=3cm, minimum height=1.2cm] at ([yshift=1.5cm]micro.center) {\small Services};
        \node[below=0.1cm of micro.south, text width=3.5cm, align=center, font=\small] {MINIX, QNX};

        % Hybrid
        \node[draw, rectangle, minimum width=3.5cm, minimum height=5cm] (hybrid) at (10,0) {};
        \node[above] at (hybrid.north) {\textbf{Hybrid}};
        \node[draw, fill=red!20, minimum width=3cm, minimum height=2.5cm] at ([yshift=-0.5cm]hybrid.center) {\small Kernel + Some};
        \node[draw, fill=blue!20, minimum width=3cm, minimum height=1.2cm] at ([yshift=1.3cm]hybrid.center) {\small Services};
        \node[below=0.1cm of hybrid.south, text width=3.5cm, align=center, font=\small] {Windows, macOS};
    \end{tikzpicture}%
    }
    \caption{Perbandingan Arsitektur Kernel}
    \label{fig:kernel-comparison}
\end{figure}

\subsection{Perbandingan Arsitektur Kernel}

\begin{table}[htbp]
    \centering
    \caption{Comparison of Kernel Architectures}
    \label{tab:kernel-comparison}
    \begin{tabular}{@{}llll@{}}
        \toprule
        \textbf{Aspek} & \textbf{Monolithic} & \textbf{Microkernel} & \textbf{Hybrid} \\
        \midrule
        Performance & Excellent & Good & Very Good \\
        Reliability & Moderate & Excellent & Good \\
        Security & Moderate & Excellent & Good \\
        Complexity & Moderate & High & High \\
        Maintainability & Moderate & Good & Moderate \\
        Examples & Linux, BSD & QNX, MINIX & Windows, macOS \\
        \bottomrule
    \end{tabular}
\end{table}

\begin{examplebox}[Tanenbaum-Torvalds Debate (1992)]
Perdebatan terkenal antara Andrew Tanenbaum (pencipta MINIX, microkernel) dan Linus Torvalds (pencipta Linux, monolithic):

\textbf{Tanenbaum}: "Linux is obsolete" - argued monolithic kernels are outdated

\textbf{Torvalds}: Defended pragmatic monolithic design - performance and simplicity

\textbf{Outcome}: Kedua pendekatan tersebut memiliki kelebihan masing-masing. Linux mendominasi servers dengan desain monolithic-modular, sementara QNX microkernel mendominasi automotive systems.
\end{examplebox}

\subsection{System Calls}

\textbf{System call} adalah mekanisme yang disediakan kernel agar program di user space dapat meminta layanan dari kernel. System calls adalah satu-satunya cara bagi aplikasi untuk berinteraksi dengan hardware melalui kernel.

\textbf{Kategori System Calls:}
\begin{enumerate}
    \item \textbf{Process Control}: \cmd{fork()}, \cmd{exec()}, \cmd{exit()}, \cmd{wait()}
    \item \textbf{File Management}: \cmd{open()}, \cmd{read()}, \cmd{write()}, \cmd{close()}
    \item \textbf{Device Management}: \cmd{ioctl()}, \cmd{read()}, \cmd{write()}
    \item \textbf{Information Maintenance}: \cmd{getpid()}, \cmd{time()}, \cmd{uname()}
    \item \textbf{Communication}: \cmd{pipe()}, \cmd{socket()}, \cmd{send()}, \cmd{receive()}
\end{enumerate}

\begin{examplebox}[System Call Flow]
Ketika program membuka file dengan \cmd{fopen()} di C:

\begin{enumerate}
    \item Application calls library function \cmd{fopen()}
    \item Library function calls system call \cmd{open()}
    \item CPU switches to kernel mode (Ring 0)
    \item Kernel validates parameters dan permissions
    \item Kernel interacts dengan file system dan disk driver
    \item Kernel returns file descriptor
    \item CPU switches back to user mode (Ring 3)
    \item Control returns to application
\end{enumerate}

Proses ini sama secara koneptual di semua OS, tetapi detil implementasi berbeda.
\end{examplebox}

%% ========================================================================
%% SESSION 2: SEJARAH & EKOSISTEM OS
%% ========================================================================

%% ------------------------------------------------------------------------
\section{Evolusi Sistem Operasi}
\label{sec:os-history}

\subsection{Timeline Singkat}

Perkembangan sistem operasi telah melalui beberapa era penting:

\begin{description}
    \item[1950s--1960s: Era Awal] Batch processing systems, program cards, no interactive computing. Contoh: IBM 701, UNIVAC.

    \item[1960s: Time-Sharing] CTSS (Compatible Time-Sharing System) di MIT - banyak pengguna berbagi waktu komputer. Proyek Multics (ambisius namun gagal, menginspirasi Unix).

    \item[1969: Lahirnya Unix] Ken Thompson dan Dennis Ritchie di Bell Labs menciptakan Unix. Sederhana, portabel (ditulis dalam C), dengan filosofi "do one thing well".

    \item[1980s: Era PC] MS-DOS (1981) untuk IBM PC. Apple Macintosh (1984) dengan GUI. Windows 1.0 (1985) sebagai shell di atas DOS.

    \item[1991: Linux] Linus Torvalds, dari Finland, merilis Linux kernel 0.01. Free, open source, Unix-compatible.

    \item[1993: Windows NT] Microsoft merilis Windows NT dengan arsitektur hybrid kernel untuk pasar enterprise.

    \item[2001: Mac OS X] Apple merilis macOS X berbasis Darwin (BSD Unix + Mach microkernel). Bersertifikasi UNIX.

    \item[2000s: Era Modern] Windows XP menyatukan lini consumer dan NT. Linux mendominasi pasar server. Momentum open source meningkat.

    \item[2007--2008: Revolusi Mobile] iOS (2007) dan Android (2008) mendefinisikan ulang komputasi dengan pendekatan mobile-first.

    \item[2010s: Era Cloud] Linux mendominasi cloud computing (AWS, Google Cloud, Azure). Container (Docker) merevolusi deployment.

    \item[2020s: Saat Ini] Komputasi hybrid, prosesor ARM (Apple Silicon), WSL membawa Linux ke Windows, pengembangan fokus keamanan.
\end{description}

%% ------------------------------------------------------------------------
\section{Ekosistem Sistem Operasi Modern}
\label{sec:os-ecosystem}

\subsection{Windows Ecosystem}

\textbf{Platform:}
\begin{itemize}
    \item \textbf{Windows 10/11} - Consumer dan Professional editions
    \item \textbf{Windows Server} - Enterprise server OS
    \item \textbf{Windows IoT} - Embedded dan industrial devices
\end{itemize}

\textbf{Karakteristik:}
\begin{itemize}
    \item Dominasi pasar desktop (pangsa pasar 75\%)
    \item Backward compatibility yang sangat baik
    \item Ekosistem perangkat lunak yang luas
    \item Integrasi enterprise (Active Directory, Group Policy)
    \item Dukungan gaming terbaik
\end{itemize}

\textbf{Teknologi:}
\begin{itemize}
    \item Filesystem: NTFS
    \item Package management: Microsoft Store, winget, Chocolatey
    \item Development: Visual Studio, .NET, PowerShell
    \item Virtualization: Hyper-V
\end{itemize}

\subsection{macOS Ecosystem}

\textbf{Karakteristik:}
\begin{itemize}
    \item Integrasi hardware-software premium (Apple Silicon)
    \item Bersertifikasi UNIX (POSIX compliant)
    \item Sangat baik untuk profesional kreatif
    \item Integrasi ekosistem iOS/iPadOS
    \item Developer-friendly (Unix tools + graphical environment)
\end{itemize}

\textbf{Teknologi:}
\begin{itemize}
    \item Filesystem: APFS (Apple File System)
    \item Package management: App Store, Homebrew
    \item Development: Xcode, Swift, Objective-C
    \item Security: Gatekeeper, FileVault, Secure Enclave
\end{itemize}

\subsection{Linux Ecosystem}

\textbf{Distribusi Linux:}

Linux bukan sistem operasi tunggal, melainkan kumpulan distribusi yang berbagi Linux kernel namun berbeda dalam package management, default software, dan target audience.

\textbf{Family Utama:}
\begin{itemize}
    \item \textbf{Debian Family} - Ubuntu, Debian, Linux Mint
    \begin{itemize}
        \item Package management: APT (dpkg)
        \item Stability-focused, large repository
        \item Popular untuk desktop dan server
    \end{itemize}

    \item \textbf{Red Hat Family} - Fedora, RHEL, CentOS Stream, Rocky Linux, AlmaLinux
    \begin{itemize}
        \item Package management: DNF/YUM (RPM)
        \item Enterprise-oriented
        \item Cutting-edge technologies
    \end{itemize}

    \item \textbf{Independent} - Arch Linux, Gentoo, Slackware
    \begin{itemize}
        \item Rolling release atau highly customizable
        \item Advanced users
        \item DIY philosophy
    \end{itemize}
\end{itemize}

\textbf{Mengapa Linux untuk Praktik?}

\begin{enumerate}
    \item \textbf{Open Source} - Gratis, source code available untuk learning
    \item \textbf{Educational} - Transparent, excellent documentation
    \item \textbf{Industry Relevance} - Dominates server, cloud, supercomputers
    \begin{itemize}
        \item 100\% of TOP500 supercomputers run Linux
        \item 90\% of public cloud workloads
        \item Most web servers
    \end{itemize}
    \item \textbf{Skills Transferable} - Unix concepts apply to macOS, BSDs
    \item \textbf{Career Value} - High demand for Linux skills in DevOps, SRE, cloud engineering
    \item \textbf{Flexibility} - Complete control, customization
    \item \textbf{Legal} - Can install unlimited copies untuk learning
\end{enumerate}

\begin{notebox}
Concepts yang dipelajari di Linux applicable ke semua Unix-like systems karena POSIX compliance. Skills seperti shell scripting, file permissions, process management work di Linux, macOS, dan BSDs.
\end{notebox}

\begin{table}[htbp]
    \centering
    \caption{Perbandingan Distribusi Linux Populer}
    \label{tab:linux-distros}
    \begin{tabular}{@{}llll@{}}
        \toprule
        \textbf{Distribusi} & \textbf{Base} & \textbf{Package Manager} & \textbf{Target User} \\
        \midrule
        Ubuntu & Debian & APT & Beginners, General \\
        Debian & Independent & APT & Stability, Servers \\
        Fedora & Red Hat & DNF & Developers, Latest Tech \\
        CentOS Stream & Red Hat & DNF & Enterprise, Testing \\
        Arch Linux & Independent & Pacman & Advanced, Customization \\
        \bottomrule
    \end{tabular}
\end{table}

%% ========================================================================
%% SESSION 3-4: INSTALASI SISTEM OPERASI
%% ========================================================================

%% ------------------------------------------------------------------------
\section{Instalasi Sistem Operasi}
\label{sec:os-installation}

\subsection{Konsep Instalasi Universal}

Instalasi sistem operasi meliputi tahapan umum berikut, berlaku untuk semua OS:

\begin{enumerate}
    \item \textbf{Persiapan} - Unduh ISO, buat media bootable, backup data
    \item \textbf{Boot} - Konfigurasi BIOS/UEFI, boot dari media instalasi
    \item \textbf{Partisi} - Alokasi ruang disk, buat partisi
    \item \textbf{Instalasi} - Salin file sistem, instal bootloader
    \item \textbf{Konfigurasi} - Buat pengguna, konfigurasi jaringan, pengaturan awal
\end{enumerate}

\subsection{BIOS vs UEFI}

\begin{table}[htbp]
    \centering
    \caption{Comparison: BIOS vs UEFI}
    \label{tab:bios-uefi}
    \begin{tabular}{@{}lll@{}}
        \toprule
        \textbf{Aspect} & \textbf{BIOS (Legacy)} & \textbf{UEFI (Modern)} \\
        \midrule
        Introduced & 1970s & 2000s \\
        Boot Mode & 16-bit & 32-bit/64-bit \\
        Partition Table & MBR & GPT \\
        Disk Limit & 2 TB & 9.4 ZB \\
        Boot Time & Slower & Faster \\
        Secure Boot & No & Yes \\
        \bottomrule
    \end{tabular}
\end{table}

\begin{warningbox}
Modern systems (post-2010) umumnya menggunakan UEFI. Secure Boot harus disabled atau configured untuk install beberapa Linux distros.
\end{warningbox}

\subsection{Metode Deployment}

\subsubsection{Virtual Machine (Recommended untuk Learning)}

\textbf{Keuntungan:}
\begin{itemize}
    \item Eksperimen yang aman - tidak mempengaruhi host system
    \item Snapshot untuk recovery yang mudah
    \item Menjalankan multiple OS secara paralel
    \item Tidak ada risiko hardware
\end{itemize}

\textbf{Software:}
\begin{itemize}
    \item VirtualBox (free, cross-platform)
    \item VMware Workstation Player (free)
    \item QEMU/KVM (Linux native, free)
    \item Hyper-V (Windows built-in)
\end{itemize}

\subsubsection{Dual Boot}

Install multiple OS pada satu machine. User memilih OS saat boot.

\textbf{Use case:} Performa native, akses hardware, penggunaan production

\textbf{Risiko:} Kesalahan partisi, konflik bootloader, kehilangan data jika tidak hati-hati

\subsubsection{Bare Metal}

Single OS installation, full hardware control.

\textbf{Use case:} Server production, workstation khusus

\subsection{Praktik: Instalasi Ubuntu Server 22.04 LTS}

\subsubsection{Prerequisites}

\begin{itemize}
    \item VirtualBox installed di host machine
    \item Ubuntu Server 22.04 LTS ISO downloaded dari \url{https://ubuntu.com/download/server}
    \item Minimum 2GB RAM, 25GB disk space untuk VM
\end{itemize}

\begin{tipbox}
Ubuntu Server dipilih karena:
\begin{itemize}
    \item LTS (Long Term Support) - 5 years updates
    \item Industry standard untuk servers
    \item Excellent documentation
    \item Large community support
\end{itemize}
\end{tipbox}

\subsubsection{Step-by-Step Installation}

\textbf{1. Create Virtual Machine}

\begin{enumerate}
    \item Buka VirtualBox, klik \keystroke{New}
    \item Name: \texttt{Ubuntu-Server-Lab}
    \item Type: Linux, Version: Ubuntu (64-bit)
    \item Memory: 2048 MB (2 GB RAM)
    \item Hard disk: Create virtual hard disk
    \begin{itemize}
        \item Type: VDI (VirtualBox Disk Image)
        \item Storage: Dynamically allocated
        \item Size: 25 GB
    \end{itemize}
\end{enumerate}

\textbf{2. Configure VM Settings}

\begin{itemize}
    \item Settings $\to$ Storage $\to$ Controller: IDE $\to$ Add optical drive
    \item Select Ubuntu Server ISO
    \item Network: NAT atau Bridged (untuk akses jaringan)
\end{itemize}

\textbf{3. Boot and Install}

\begin{enumerate}
    \item Start VM, akan boot dari ISO
    \item Select language: English
    \item Keyboard configuration: sesuai keyboard Anda
    \item Network configuration:
    \begin{itemize}
        \item DHCP (automatic) - direkomendasikan untuk pemula
        \item Atau konfigurasi static IP (advanced)
    \end{itemize}
    \item Proxy: Kosongkan (kecuali diperlukan)
    \item Mirror: Default Ubuntu archive mirror
\end{enumerate}

\textbf{4. Partitioning}

Pilihan partitioning:

\textbf{Option 1: Guided - use entire disk (Direkomendasikan untuk pemula)}
\begin{itemize}
    \item Automatic partitioning
    \item Mengatur /, /boot/efi, dan swap
    \item Sederhana dan cepat
\end{itemize}

\textbf{Option 2: Custom storage layout (Advanced)}
\begin{itemize}
    \item Manual partition creation
    \item Dijelaskan detail di section~\ref{sec:disk-partitioning}
\end{itemize}

\textbf{5. Profile Setup}

\begin{lstlisting}[language=bash, caption={User account information}]
Your name: John Doe
Your server's name: ubuntu-server
Pick a username: john
Choose a password: [strong password]
Confirm your password: [repeat password]
\end{lstlisting}

\textbf{6. SSH Setup}

\begin{itemize}
    \item Install OpenSSH server: \textbf{YES} (highly recommended)
    \item Import SSH identity: No (skip untuk now)
\end{itemize}

\begin{notebox}
OpenSSH server memungkinkan remote access ke server via terminal. Essential untuk production servers.
\end{notebox}

\textbf{7. Featured Server Snaps}

Optional packages:
\begin{itemize}
    \item Docker - container runtime
    \item Microk8s - Kubernetes
    \item (Skip untuk basic installation)
\end{itemize}

\textbf{8. Installation Progress}

\begin{itemize}
    \item Instalasi akan mengunduh dan memasang paket (~15-20 menit)
    \item Tunggu hingga "Installation complete!" muncul
    \item Pilih "Reboot Now"
    \item Keluarkan media instalasi (VirtualBox otomatis mengeluarkannya)
\end{itemize}

\textbf{9. First Login}

\begin{lstlisting}[language=bash, caption={Login dan verify installation}]
# Login dengan username yang dibuat
ubuntu-server login: john
Password: [enter password]

# Verify system information
$ uname -r
5.15.0-56-generic

$ cat /etc/os-release
NAME="Ubuntu"
VERSION="22.04.1 LTS (Jammy Jellyfish)"

$ ip a
# Should show network interfaces dengan IP address
\end{lstlisting}

\begin{examplebox}[Snapshot VM]
Setelah instalasi berhasil, buat snapshot di VirtualBox:
\begin{enumerate}
    \item Matikan VM
    \item VirtualBox $\to$ VM $\to$ Snapshots $\to$ Take
    \item Name: "Fresh Install"
\end{enumerate}

Ini memungkinkan restore ke kondisi bersih kapanpun dibutuhkan.
\end{examplebox}

%% ========================================================================
%% SESSION 4-5: PARTISI DAN FILESYSTEM
%% ========================================================================

%% ------------------------------------------------------------------------
\section{Partisi Disk dan Sistem File}
\label{sec:disk-partitioning}

\subsection{Konsep Partisi}

\textbf{Partisi} adalah logical division dari physical disk. Setiap partisi diperlakukan sebagai independent storage unit oleh OS.

\textbf{Mengapa perlu partisi?}
\begin{itemize}
    \item \textbf{Organisasi} - Memisahkan OS, data pengguna, aplikasi
    \item \textbf{Multi-OS} - Instalasi multiple OS pada satu disk
    \item \textbf{Backup/Recovery} - Lebih mudah untuk backup partisi tertentu
    \item \textbf{Keamanan} - Mengisolasi data dengan permission yang berbeda
    \item \textbf{Performa} - Optimasi untuk jenis workload yang berbeda
\end{itemize}

\subsection{Partition Tables}

\subsubsection{MBR (Master Boot Record)}

\textbf{Legacy scheme} - digunakan dengan BIOS:
\begin{itemize}
    \item Maksimum 4 primary partitions
    \item Extended partition dapat berisi multiple logical partitions
    \item Ukuran disk maksimum: 2 TB
    \item Ukuran partition table: 64 bytes
\end{itemize}

\subsubsection{GPT (GUID Partition Table)}

\textbf{Modern scheme} - digunakan dengan UEFI:
\begin{itemize}
    \item Praktis unlimited partitions (128 di Windows/Linux)
    \item Ukuran disk maksimum: 9.4 ZB (zettabytes)
    \item Redundant partition tables untuk reliability
    \item Diperlukan untuk UEFI Secure Boot
    \item Backward compatible dengan BIOS (Protective MBR)
\end{itemize}

\subsection{Skema Partisi Linux}

\subsubsection{Minimum Partitioning}

\begin{lstlisting}[caption={Minimal partition scheme}]
/           20 GB     ext4      Root filesystem
swap        2-4 GB    swap      Virtual memory
\end{lstlisting}

\subsubsection{Recommended Partitioning}

\begin{lstlisting}[caption={Recommended partition scheme}]
/boot/efi   512 MB    FAT32     EFI System Partition (UEFI)
/boot       1 GB      ext4      Kernel dan bootloader
/           20 GB     ext4      Root filesystem
/home       50+ GB    ext4      User data
swap        4 GB      swap      Virtual memory
\end{lstlisting}

\textbf{Penjelasan Mount Points:}

\begin{description}
    \item[\texttt{/}] (root) Sistem files, installed programs, configurations
    \item[\texttt{/boot}] Kernel, initramfs, GRUB bootloader
    \item[\texttt{/boot/efi}] EFI bootloader (UEFI systems only)
    \item[\texttt{/home}] User directories dan personal data
    \item[\texttt{swap}] Virtual memory, hibernation support
\end{description}

\begin{tipbox}
\textbf{Mengapa memisahkan \texttt{/home}?}

Keuntungan:
\begin{itemize}
    \item Reinstall OS tanpa kehilangan data pengguna
    \item Backup lebih mudah (backup hanya \texttt{/home})
    \item Upgrade distro dengan mempertahankan pengaturan
    \item Filesystem berbeda atau enkripsi untuk data
\end{itemize}
\end{tipbox}

\subsection{Swap Space}

\textbf{Swap} adalah disk space used sebagai extension dari RAM (virtual memory).

\textbf{Kapan swap digunakan?}
\begin{itemize}
    \item Saat RAM penuh, inactive pages dipindahkan ke swap
    \item Hibernation (suspend-to-disk) - menyimpan konten RAM ke swap
\end{itemize}

\textbf{Swap Sizing Guidelines:}

\begin{table}[htbp]
    \centering
    \caption{Rekomendasi Ukuran Swap}
    \label{tab:swap-sizing}
    \begin{tabular}{@{}ll@{}}
        \toprule
        \textbf{RAM Size} & \textbf{Swap Size} \\
        \midrule
        $\leq$ 2 GB & 2x RAM \\
        2--8 GB & Same as RAM \\
        $>$ 8 GB & 4--8 GB fixed \\
        Hibernation needed & RAM size + 2 GB \\
        \bottomrule
    \end{tabular}
\end{table}

\begin{warningbox}
  Swap bukanlah pengganti RAM yang tidak mencukupi. Menambahkan RAM jauh lebih efektif daripada membuat swap yang besar.
\end{warningbox}

\subsection{Filesystem Types}

\subsubsection{Cross-Platform Comparison}

\begin{table}[htbp]
    \centering
    \caption{Perbandingan Filesystem}
    \label{tab:filesystem-comparison}
    \begin{tabular}{@{}lllll@{}}
        \toprule
        \textbf{Filesystem} & \textbf{Platform} & \textbf{Max File} & \textbf{Journaling} & \textbf{Use Case} \\
        \midrule
        NTFS & Windows & 16 TB & Yes & Windows system \\
        APFS & macOS & 8 EB & CoW & Modern macOS \\
        ext4 & Linux & 16 TB & Yes & Linux default \\
        XFS & Linux & 8 EB & Yes & Large files \\
        Btrfs & Linux & 16 EB & CoW & Advanced features \\
        exFAT & Universal & 16 EB & No & USB drives \\
        \bottomrule
    \end{tabular}
\end{table}

\subsubsection{Linux Filesystems}

\textbf{ext4 (Fourth Extended Filesystem)}
\begin{itemize}
    \item Default untuk sebagian besar distribusi Linux
    \item Matang, stabil, dan teruji dengan baik
    \item Journaling untuk crash recovery
    \item Performa baik untuk penggunaan umum
\end{itemize}

\noindent
\textbf{XFS}
\begin{itemize}
    \item Filesystem berkinerja tinggi
    \item Sangat baik untuk large files dan parallel I/O
    \item Default di Red Hat Enterprise Linux
    \item Tidak dapat diperkecil (hanya diperbesar)
\end{itemize}

\noindent
\textbf{Btrfs (B-tree filesystem)}
\begin{itemize}
    \item Filesystem copy-on-write modern
    \item Snapshot dan kompresi built-in
    \item Dukungan RAID
    \item Masih dalam pematangan (production-ready di openSUSE)
\end{itemize}

\subsection{Praktik: Manual Partitioning}

Untuk pengguna advanced, manual partitioning memberikan kontrol penuh.

Selama instalasi Ubuntu, pilih "Custom storage layout":

\begin{lstlisting}[caption={Example manual partition layout}]
Device      Size    Type    Mount Point    Filesystem
/dev/sda1   512 MB  EFI     /boot/efi      FAT32
/dev/sda2   1 GB    ext4    /boot          ext4
/dev/sda3   20 GB   ext4    /              ext4
/dev/sda4   50 GB   ext4    /home          ext4
/dev/sda5   4 GB    swap    [swap]         swap
\end{lstlisting}

\subsubsection{Post-Installation Filesystem Commands}

\begin{lstlisting}[language=bash, caption={Filesystem verification dan management}]
# Check disk usage
df -h

# Check disk partitions
lsblk
sudo fdisk -l

# Check filesystem type
lsblk -f
blkid

# Mount additional filesystem manually
sudo mkdir /mnt/data
sudo mount /dev/sdb1 /mnt/data

# Permanent mount - edit /etc/fstab
cat /etc/fstab
# Format: <device> <mount point> <type> <options> <dump> <pass>

# UUID-based mount (preferred)
UUID=xxx-xxx-xxx  /data  ext4  defaults  0  2
\end{lstlisting}

\begin{notebox}
\textbf{UUID vs Device Names}

Gunakan UUID daripada \texttt{/dev/sdX} di fstab karena:
\begin{itemize}
    \item Device names dapat berubah (sda $\to$ sdb jika menambah disk)
    \item UUID adalah persistent identifier
    \item Lebih reliable untuk production systems
\end{itemize}

Dapatkan UUID: \cmd{blkid} atau \cmd{lsblk -f}
\end{notebox}

%% ========================================================================
%% SESSION 5-6: POST-INSTALLATION & SUMMARY
%% ========================================================================

%% ------------------------------------------------------------------------
\section{Konfigurasi Post-Installation}
\label{sec:post-installation}

Setelah instalasi selesai, beberapa konfigurasi penting harus dilakukan.

\subsection{System Updates}

\textbf{Mengapa penting?}
\begin{itemize}
    \item Security patches untuk vulnerabilities
    \item Bug fixes dan stability improvements
    \item New features dan driver updates
\end{itemize}

\begin{lstlisting}[language=bash, caption={Update sistem Ubuntu/Debian}]
# Update package lists dari repositories
sudo apt update

# Upgrade paket yang terinstal ke versi terbaru
sudo apt upgrade

# Full system upgrade (menangani dependencies)
sudo apt full-upgrade

# Hapus paket yang tidak dibutuhkan
sudo apt autoremove

# Bersihkan package cache
sudo apt clean
\end{lstlisting}

\subsection{Software Installation Essentials}

\begin{lstlisting}[language=bash, caption={Install essential development tools}]
# Development tools
sudo apt install -y build-essential

# Version control
sudo apt install -y git

# Network utilities
sudo apt install -y curl wget net-tools

# Text editors
sudo apt install -y vim nano

# System monitoring
sudo apt install -y htop iotop

# Complete essentials in one command
sudo apt install -y \
    build-essential \
    git \
    curl \
    wget \
    vim \
    htop \
    tree \
    net-tools \
    openssh-server
\end{lstlisting}

\subsection{Verifikasi Koneksi Jaringan}

Pastikan sistem dapat terhubung ke internet.

\begin{lstlisting}[language=bash, caption={Check network status}]
# Check IP address
ip addr show

# Test internet connectivity
ping -c 4 google.com

# Check DNS
cat /etc/resolv.conf
\end{lstlisting}

\begin{tipbox}
Jika menggunakan VirtualBox, pastikan network adapter di-set ke "NAT" atau "Bridged Adapter" untuk koneksi internet.
\end{tipbox}

\subsection{Akses SSH (Opsional)}

Jika ingin remote access ke server melalui SSH (jika sudah terinstall saat instalasi):

\begin{lstlisting}[language=bash, caption={Check SSH status}]
# Check if SSH installed dan running
sudo systemctl status ssh

# Get IP address untuk remote access
ip addr show
\end{lstlisting}

\begin{examplebox}[Remote Login via SSH]
Dari komputer lain pada jaringan yang sama:
\begin{lstlisting}[language=bash]
ssh username@192.168.1.100
\end{lstlisting}

Ganti alamat IP dengan alamat IP server.
\end{examplebox}

\subsection{Menampilkan Informasi Sistem}

\subsubsection{Install Neofetch}

Neofetch adalah tool populer untuk menampilkan informasi sistem dengan visual yang menarik.

\begin{lstlisting}[language=bash, caption={Install neofetch}]
# Install neofetch
sudo apt install -y neofetch

# Run neofetch
neofetch
\end{lstlisting}

Neofetch akan menampilkan:
\begin{itemize}
    \item Logo distro Linux
    \item Informasi OS (distro, versi, kernel)
    \item CPU dan memory usage
    \item Disk space
    \item Uptime sistem
    \item Dan informasi lainnya dalam format yang mudah dibaca
\end{itemize}

\begin{tipbox}
Add \cmd{neofetch} ke file \file{\textasciitilde/.bashrc} agar otomatis tampil setiap login:
\begin{lstlisting}[language=bash]
echo "neofetch" >> ~/.bashrc
\end{lstlisting}
\end{tipbox}

\subsubsection{Perintah Informasi Sistem Dasar}

\begin{lstlisting}[language=bash, caption={Perintah informasi sistem dasar}]
# Informasi OS
cat /etc/os-release
hostnamectl

# Versi kernel
uname -r

# Informasi CPU
lscpu

# Informasi memory
free -h

# Informasi disk
df -h
lsblk

# System uptime
uptime

# Who is logged in
who
\end{lstlisting}

\begin{examplebox}[Praktik: Explore System Info]
Coba jalankan perintah berikut dan perhatikan output:
\begin{enumerate}
    \item \cmd{neofetch} -- Tampilan lengkap dengan visual
    \item \cmd{free -h} -- Check RAM usage
    \item \cmd{df -h} -- Check disk space
    \item \cmd{lscpu} -- Detail CPU
    \item \cmd{uptime} -- Berapa lama sistem running
\end{enumerate}
\end{examplebox}

%% ------------------------------------------------------------------------
\section{Rangkuman}
\label{sec:summary}

Dalam bab ini, kita telah mempelajari konsep fundamental sistem operasi dan praktik instalasi:

\subsection{Konsep Sistem Operasi}

\begin{itemize}
    \item \textbf{Definisi}: OS adalah resource manager dan extended machine
    \item \textbf{Fungsi Utama}: Process, Memory, File, I/O Management, dan Security
    \item \textbf{Jenis OS}: Desktop, Server, Mobile, Real-Time, Embedded
    \item \textbf{Kernel Architectures}:
    \begin{itemize}
        \item Monolithic (Linux) - Performance
        \item Microkernel (QNX) - Reliability
        \item Hybrid (Windows, macOS) - Balance
    \end{itemize}
    \item \textbf{System Calls}: Interface between user space dan kernel services
\end{itemize}

\subsection{Perbandingan Sistem Operasi}

\begin{itemize}
    \item \textbf{Windows}: Desktop dominance, gaming, business software
    \item \textbf{macOS}: Creative professionals, hardware integration
    \item \textbf{Linux}: Server dominance, open source, flexibility
    \item \textbf{Android/iOS}: Mobile platforms
\end{itemize}

\subsection{Practical Skills}

\begin{itemize}
    \item Installation process (VM setup, Ubuntu Server)
    \item Partitioning concepts (MBR vs GPT, partition schemes)
    \item Filesystem understanding (ext4, swap, mount points)
    \item Post-installation: system updates dan basic configuration
    \item System information gathering dengan neofetch dan basic commands
\end{itemize}

\subsection{Key Takeaways}

\begin{enumerate}
    \item Konsep OS bersifat universal - berlaku di berbagai platform
    \item Linux sangat baik untuk pembelajaran: terbuka, gratis, relevan dengan industri
    \item Virtual machines menyediakan lingkungan pembelajaran yang aman
    \item Partisi yang tepat penting untuk fleksibilitas dan recovery
    \item System updates harus dilakukan secara berkala
    \item Basic system monitoring tools membantu memahami kondisi sistem
\end{enumerate}

%% ------------------------------------------------------------------------
\section{Latihan}
\label{sec:exercises}

\subsection{Latihan Konseptual}

\begin{exercisebox}[1.1]
Jelaskan 5 fungsi utama sistem operasi dengan contoh konkret dari minimal 2 OS berbeda (Windows, macOS, atau Linux).
\end{exercisebox}

\begin{exercisebox}[1.2]
Kapan sebaiknya menggunakan Windows vs Linux vs macOS? Analisis berdasarkan use case: gaming, development, server, creative work, dan enterprise.
\end{exercisebox}

\subsection{Latihan Praktikal}

\begin{exercisebox}[1.3]
Install Ubuntu Server 22.04 LTS di VirtualBox dengan langkah berikut:
\begin{enumerate}
    \item Unduh Ubuntu Server ISO dari website resmi
    \item Buat VM baru di VirtualBox (RAM: 2GB, Disk: 25GB)
    \item Install dengan automatic partitioning (guided)
    \item Buat user account dengan password yang kuat
    \item Reboot dan login ke sistem
    \item Dokumentasikan proses instalasi dengan screenshot langkah-langkah kunci
\end{enumerate}
\end{exercisebox}

\begin{exercisebox}[1.4]
Setelah instalasi Ubuntu Server, lakukan tasks berikut:
\begin{enumerate}
    \item Update package list: \cmd{sudo apt update}
    \item Upgrade packages: \cmd{sudo apt upgrade}
    \item Install neofetch: \cmd{sudo apt install neofetch}
    \item Jalankan neofetch dan screenshot hasilnya
    \item Check disk usage dengan \cmd{df -h}
    \item Check memory dengan \cmd{free -h}
    \item Dokumentasikan output dari setiap command
\end{enumerate}
\end{exercisebox}

\begin{exercisebox}[1.5]
Eksplorasi sistem yang baru diinstall:
\begin{enumerate}
    \item Tampilkan informasi OS: \cmd{cat /etc/os-release}
    \item Tampilkan versi kernel: \cmd{uname -r}
    \item List partisi: \cmd{lsblk}
    \item Check network connectivity: \cmd{ping -c 4 google.com}
    \item Install dan jalankan \cmd{htop} untuk melihat resource usage
    \item Buat laporan singkat tentang konfigurasi sistem Anda
\end{enumerate}
\end{exercisebox}

\subsection{Latihan Refleksi}

\begin{exercisebox}[1.6]
Ceritakan pengalaman Anda dengan sistem operasi:
\begin{enumerate}
    \item Sistem operasi apa yang Anda gunakan sehari-hari? (Windows, macOS, Linux, atau lainnya)
    \item Berapa lama Anda menggunakan sistem operasi tersebut?
    \item Apa yang Anda sukai dari sistem operasi tersebut?
    \item Apa tantangan atau masalah yang pernah Anda hadapi?
    \item Apakah Anda pernah menggunakan sistem operasi lain? Bandingkan pengalaman Anda.
    \item Setelah mempelajari bab ini, apakah ada sistem operasi lain yang ingin Anda coba? Mengapa?
\end{enumerate}
Tulis refleksi Anda dalam 300-500 kata disertai dengan dokumentasi.
\end{exercisebox}

%% ------------------------------------------------------------------------
\section{Referensi}
\label{sec:references}

\subsection{Buku Rekomendasi}

\begin{itemize}
    \item \textbf{Operating System Concepts (10th Ed)} - Silberschatz, Galvin, Gagne \\
    Classic textbook untuk OS theory

    \item \textbf{Modern Operating Systems (4th Ed)} - Andrew Tanenbaum \\
    Comprehensive OS concepts dengan implementation details

    \item \textbf{How Linux Works (3rd Ed)} - Brian Ward \\
    Practical guide untuk understanding Linux internals

    \item \textbf{The Linux Command Line (2nd Ed)} - William Shotts \\
    Excellent resource untuk command-line mastery

    \item \textbf{UNIX and Linux System Administration Handbook (5th Ed)} - Nemeth et al. \\
    Industry standard untuk system administration
\end{itemize}

\subsection{Online Resources}

\textbf{Documentation:}
\begin{itemize}
    \item Linux Foundation: \url{https://www.linuxfoundation.org}
    \item Ubuntu Documentation: \url{https://ubuntu.com/server/docs}
    \item Arch Wiki: \url{https://wiki.archlinux.org} (excellent cross-distro reference)
    \item Red Hat Documentation: \url{https://access.redhat.com/documentation}
\end{itemize}

\textbf{Interactive Learning:}
\begin{itemize}
    \item Linux Journey: \url{https://linuxjourney.com}
    \item OverTheWire Bandit: \url{https://overthewire.org/wargames/bandit/}
    \item edX: Introduction to Linux (LFS101x)
\end{itemize}

\textbf{Official Sites:}
\begin{itemize}
    \item The Linux Kernel: \url{https://www.kernel.org}
    \item Microsoft Docs: \url{https://docs.microsoft.com}
    \item Apple Developer: \url{https://developer.apple.com/documentation}
\end{itemize}
